\section{Acyclicity criteria by coverings}
\label{section-covering-acyclicity-criteria}

\noindent
We first extend the notation for the \v Cech complex to the coverings by
subsets which occur below. Let $A$ be a topological space, let
$\mathcal R=(R_i)_{i\in I}$ be an indexed covering of $A$ by subsets, and
let $\mathcal F$ be an abelian sheaf on $A$. Put
$$
R_{i_0\ldots i_p}=R_{i_0}\cap\ldots\cap R_{i_p}
$$
and define
\begin{align*}
C^p(\mathcal R,\underline H^q(\mathcal F))
&=
\prod_{(i_0,\ldots,i_p)\in I^{p+1}}
H^q(R_{i_0\ldots i_p},
\mathcal F|_{R_{i_0\ldots i_p}}),\\
H^p(\mathcal R,\underline H^q(\mathcal F))
&=
H^p(C^\bullet(\mathcal R,\underline H^q(\mathcal F))).
\end{align*}
The differential is the alternating sum of the restriction maps, as in
Equation (\ref{equation-d-cech}). For $q=0$ we abbreviate these groups to
$C^p(\mathcal R,\mathcal F)$ and $H^p(\mathcal R,\mathcal F)$.
If $A\subset X$ and $\mathcal F$ is a sheaf on $X$, this notation always
means that the construction is applied to $\mathcal F|_A$.

If $L$ is a finite collection of coverings of $A$, then
$\mathcal R^L$ denotes their common intersection covering: its members are
the nonempty intersections obtained by choosing one member of every
covering in $L$. We use the convention that $\mathcal R^\emptyset$ is the
one-member covering $(A)$. The nerve of $\mathcal R$ has dimension at most
$d$ if no $d+2$ distinct members of $\mathcal R$ have nonempty
intersection.

\begin{lemma}[Leray--Cartan spectral sequence for closed coverings]
\label{lemma-cech-spectral-sequence-closed-covering}
\begin{reference}
\cite[Chapter III, \S 3.8, Theorem 3.8.1 and the following
remark]{Tohoku}
\end{reference}
Let $A$ be a topological space and let
$\mathcal R=(R_i)_{i\in I}$ be a covering. Assume one of the following.
\begin{enumerate}
\item Every $R_i$ is open in $A$.
\item The set $I$ is finite and every $R_i$ is closed in $A$.
\item The space $A$ is paracompact and $(R_i)_{i\in I}$ is a locally
finite covering by closed subsets.
\end{enumerate}
For every abelian sheaf $\mathcal F$ on $A$ there is a functorial first
quadrant spectral sequence
$$
E_1^{p,q}=
\prod_{(i_0,\ldots,i_p)\in I^{p+1}}
H^q(R_{i_0\ldots i_p},
\mathcal F|_{R_{i_0\ldots i_p}}),
\qquad
E_2^{p,q}=H^p(\mathcal R,\underline H^q(\mathcal F))
\Longrightarrow H^{p+q}(A,\mathcal F).
$$
Its bottom edge maps are canonical maps
$$
H^n(\mathcal R,\mathcal F)\longrightarrow H^n(A,\mathcal F),
$$
and its other edge maps are the restriction maps
$$
H^n(A,\mathcal F)\longrightarrow
H^0(\mathcal R,\underline H^n(\mathcal F)).
$$

These spectral sequences and both edge maps are natural for refinement of
coverings and for trace on a subset, whenever the source and target
coverings satisfy one of the three conditions above. A refinement function
$c:J\to I$ with $V_j\subset R_{c(j)}$ induces the usual restriction map,
independently of $c$ from the $E_2$-page onward, and induces the identity on
$H^*(A,\mathcal F)$. If $B\subset A$, trace induces the restriction map
$H^*(A,\mathcal F)\to H^*(B,\mathcal F|_B)$ on the abutments.

If $n>0$ and
$$
H^q(R_{i_0\ldots i_p},
\mathcal F|_{R_{i_0\ldots i_p}})=0
\quad\text{for }0<q<n,
$$
for all $p$ and all tuples of indices, then the edge maps give an exact
sequence
$$
H^n(\mathcal R,\mathcal F)\longrightarrow H^n(A,\mathcal F)
\longrightarrow H^0(\mathcal R,\underline H^n(\mathcal F)).
$$
\end{lemma}

\begin{proof}
In the open case this is Lemma \ref{lemma-cech-spectral-sequence}, using the
covering resolution of Lemma \ref{lemma-covering-resolution}. We give the
closed-cover argument because it is needed below.

For a tuple $\mathbf i=(i_0,\ldots,i_p)$, let
$j_{\mathbf i}:R_{\mathbf i}\to A$ be the closed inclusion. Define
$$
\mathcal K^p(\mathcal R,\mathcal F)=
\prod_{\mathbf i\in I^{p+1}}
(j_{\mathbf i})_*(\mathcal F|_{R_{\mathbf i}})
$$
with the alternating restriction differential. There is an augmentation
$$
\mathcal F\longrightarrow
\mathcal K^\bullet(\mathcal R,\mathcal F).
$$
This is a quasi-isomorphism. Indeed, in the finite case, and locally in the
locally finite case, stalks commute with the displayed products. At a point
$x$, the nonzero factors form the augmented cochain complex of the simplex
on the nonempty set
$$
I_x=\{i\in I\mid x\in R_i\}
$$
with coefficients in $\mathcal F_x$. Choosing one element of $I_x$ gives
the usual contracting homotopy.

We recall a convenient construction of the hypercohomology spectral
sequence of this resolution. Apply the exact Godement resolution functor of
Lemma \ref{lemma-godement-resolution} to every term
$\mathcal K^p(\mathcal R,\mathcal F)$. For each Godement degree its
exactness turns the augmented covering resolution above into a flasque
resolution of the corresponding Godement term of $\mathcal F$. Hence one
spectral sequence of the resulting double complex of global sections
collapses to the complex of global sections of the Godement resolution of
$\mathcal F$, and its total cohomology is $H^*(A,\mathcal F)$. The other
spectral sequence has
$$
E_1^{p,q}=H^q(A,\mathcal K^p(\mathcal R,\mathcal F)).
$$

Suppose first that $I$ is finite. Direct image by a closed inclusion is
exact, it sends injectives to injectives, and its global sections are the
global sections on the closed subset. Finite products commute with
cohomology. Consequently
$$
H^q(A,\mathcal K^p(\mathcal R,\mathcal F))=
\prod_{\mathbf i\in I^{p+1}}
H^q(R_{\mathbf i},\mathcal F|_{R_{\mathbf i}}).
$$

Suppose next that $A$ is paracompact and the closed covering is locally
finite. Choose a flasque resolution
$\mathcal F\to\mathcal I^\bullet$. Every $\mathcal I^r$ is soft by
Lemma \ref{lemma-flasque-soft-family-supports}, and its restriction to a
closed subset is soft. For fixed $p$, the family of closed subsets
$R_{\mathbf i}$, indexed by $I^{p+1}$, is locally finite. It follows
directly from the extension property that
$\mathcal K^p(\mathcal R,\mathcal I^r)$ is soft: a section over a closed
subset is a locally finite family of sections over its intersections with
the $R_{\mathbf i}$, and the members of this family can be extended one by
one. These sheaves are acyclic by Lemma
\ref{lemma-family-supports-flasque-soft-acyclic}.

The functor $\mathcal K^p(\mathcal R,-)$ is exact in this case, as is seen
on stalks where only finitely many factors occur. Thus
$\mathcal K^p(\mathcal R,\mathcal I^\bullet)$ computes the cohomology of
$\mathcal K^p(\mathcal R,\mathcal F)$. Each closed subspace
$R_{\mathbf i}$ is paracompact, and
$\mathcal I^\bullet|_{R_{\mathbf i}}$ is a soft resolution. Since products
of abelian groups are exact, we again obtain
$$
H^q(A,\mathcal K^p(\mathcal R,\mathcal F))=
\prod_{\mathbf i\in I^{p+1}}
H^q(R_{\mathbf i},\mathcal F|_{R_{\mathbf i}}).
$$
This proves the displayed spectral sequence.

The augmentation identifies the bottom and vertical edge maps with the
maps in the statement. In particular, the vertical edge sends a class to
the compatible family of its restrictions to the $R_i$. A refinement
function gives a morphism of the covering resolutions by restriction.
Two refinement functions give homotopic morphisms by the usual prism
homotopy, so the induced maps agree from $E_2$ onward. The augmentation
shows that the map on the common abutment is the identity. Restricting all
terms to $B$ gives the trace map and the restriction map on the abutments.
This also proves all asserted naturality.

Finally, under the last vanishing hypothesis the terms with
$0<q<n$ cannot contribute in total degree $n$. Thus the filtration in that
degree gives the asserted exact edge sequence.
\end{proof}

\begin{lemma}[Coverwise local killing]
\label{lemma-covering-acyclicity-local-killing}
\begin{reference}
\cite[Chapter III, \S 3.9, Theorem 3.9.1, pp. 179--180]{Tohoku}
\end{reference}
Let $X$ be a topological space and let $\mathcal S$ be a nonempty
collection of subsets of $X$. For every $A\in\mathcal S$, let
$\mathfrak R(A)$ be a nonempty collection of indexed coverings of $A$.
Assume the following.
\begin{enumerate}
\item The members of every $\mathcal R\in\mathfrak R(A)$ and all their
nonempty finite intersections belong to $\mathcal S$.
\item If $\mathcal R,\mathcal R'\in\mathfrak R(A)$ and $B$ is a member of
$\mathcal R$, then the trace $\mathcal R'|_B$ belongs to
$\mathfrak R(B)$.
\item One of the following three alternatives holds: all members of
$\mathcal S$ are open in $X$; all members are closed and all coverings in
the collections $\mathfrak R(A)$ are finite; or all members are closed,
$X$ is paracompact, and all those coverings are locally finite.
\end{enumerate}
Let $\mathcal F$ be an abelian sheaf on $X$ and let $n\geq0$. Suppose
\begin{enumerate}
\item[$A(n)$] $H^i(\mathcal R,\mathcal F)=0$ for
$1\leq i\leq n$, every $A\in\mathcal S$, and every
$\mathcal R\in\mathfrak R(A)$.
\item[$B(n-1)$] For $A\in\mathcal S$, $1\leq i\leq n-1$, and
$c\in H^i(A,\mathcal F|_A)$, there is a finite
$L\subset\mathfrak R(A)$ such that $c$ restricts to zero on every member
of $\mathcal R^L$.
\end{enumerate}
Then
$$
H^i(A,\mathcal F|_A)=0
\quad\text{for }A\in\mathcal S\text{ and }1\leq i\leq n-1.
$$
Moreover, a class $c\in H^n(A,\mathcal F|_A)$ is zero if and only if
there is a finite $L\subset\mathfrak R(A)$ such that $c$ restricts to
zero on every member of $\mathcal R^L$.
\end{lemma}

\begin{proof}
We argue by induction on $n$. The assertion for $n=0$ follows from the
convention $\mathcal R^\emptyset=(A)$, or directly by checking germs.
Suppose $n>0$. Applying the induction hypothesis in degree $n-1$ and using
$B(n-1)$ shows that
$$
H^i(A,\mathcal F|_A)=0
\quad\text{for }1\leq i\leq n-1
$$
and every $A\in\mathcal S$.

It remains to prove the characterization in degree $n$. The forward
implication follows by taking $L=\emptyset$. For the converse, suppose
$c\in H^n(A,\mathcal F|_A)$ is killed on the members of
$\mathcal R^L$. We use induction on $|L|$. The case $|L|=0$ is immediate.
If $L=\{\mathcal R\}$, all the nonzero intersections of members of
$\mathcal R$ belong to $\mathcal S$, so their cohomology vanishes in
degrees $1,\ldots,n-1$. Lemma
\ref{lemma-cech-spectral-sequence-closed-covering} therefore gives the exact
sequence
$$
H^n(\mathcal R,\mathcal F)\longrightarrow
H^n(A,\mathcal F|_A)\longrightarrow
H^0(\mathcal R,\underline H^n(\mathcal F)).
$$
The first group is zero by $A(n)$, and the image of $c$ in the last group
is its family of restrictions to the members of $\mathcal R$. Hence $c=0$.

Now let $L=\{\mathcal R^1,\ldots,\mathcal R^k\}$ with $k>1$. For a member
$B$ of $\mathcal R^1$, the trace coverings
$\mathcal R^2|_B,\ldots,\mathcal R^k|_B$ belong to $\mathfrak R(B)$.
Trace naturality in the preceding lemma says that $c|_B$ is killed on their
common intersection covering. The induction on $k$ gives $c|_B=0$ for
every member $B$ of $\mathcal R^1$. The case $|L|=1$ now gives $c=0$.
\end{proof}

\begin{lemma}[Equivalent covering acyclicity conditions]
\label{lemma-covering-acyclicity-conditions}
\begin{reference}
\cite[Chapter III, \S 3.9, Corollary 1 to Theorem 3.9.1,
p. 179]{Tohoku}
\end{reference}
In the situation of Lemma
\ref{lemma-covering-acyclicity-local-killing}, the following are equivalent.
\begin{enumerate}
\item $H^i(A,\mathcal F|_A)=0$ for every $A\in\mathcal S$ and
$1\leq i\leq n$.
\item Conditions $A(n)$ and $B(n)$ hold.
\end{enumerate}
\end{lemma}

\begin{proof}
Assume (1). Then $B(n)$ is immediate. In the spectral sequence of Lemma
\ref{lemma-cech-spectral-sequence-closed-covering}, every finite
intersection of members of a covering belongs to $\mathcal S$. Thus, in
total degrees at most $n$, only the row $q=0$ can be nonzero, and the bottom
edge maps identify
$$
H^i(\mathcal R,\mathcal F)=H^i(A,\mathcal F|_A)
\quad\text{for }0\leq i\leq n.
$$
This proves $A(n)$.

Conversely, apply Lemma \ref{lemma-covering-acyclicity-local-killing} using
$A(n)$ and $B(n-1)$. It gives the vanishing through degree $n-1$ and says
that a degree $n$ class is zero exactly when it is killed on a common
intersection covering. Condition $B(n)$ kills every such class.
\end{proof}

\begin{lemma}[Special coverings compute acyclicity]
\label{lemma-special-coverings-compute-acyclicity}
\begin{reference}
\cite[Chapter III, \S 3.9, Corollary 2 to Theorem 3.9.1,
p. 179]{Tohoku}
\end{reference}
In the situation of Lemma
\ref{lemma-covering-acyclicity-local-killing}, assume in addition that for
every $A\in\mathcal S$ and every open covering $\mathcal U$ of $A$, there
is a finite $L\subset\mathfrak R(A)$ such that $\mathcal R^L$ refines
$\mathcal U$. Then the following are equivalent.
\begin{enumerate}
\item $H^i(A,\mathcal F|_A)=0$ for every $A\in\mathcal S$ and
$1\leq i\leq n$.
\item $H^i(\mathcal R,\mathcal F)=0$ for every $A\in\mathcal S$,
$\mathcal R\in\mathfrak R(A)$, and $1\leq i\leq n$.
\end{enumerate}
\end{lemma}

\begin{proof}
Lemma \ref{lemma-kill-cohomology-class-on-covering}, applied on the
ringed space $(A,\underline{\mathbf Z}_A)$, gives an open covering on which
any positive-degree class becomes zero. By hypothesis a common intersection
covering $\mathcal R^L$ refines that open covering. Thus $B(m)$ holds for
every $m$. The result is now Lemma
\ref{lemma-covering-acyclicity-conditions}.
\end{proof}

\begin{lemma}[Bounded nerves and acyclicity]
\label{lemma-bounded-nerve-special-coverings-acyclic}
\begin{reference}
\cite[Chapter III, \S 3.9, Corollary 3 to Theorem 3.9.1,
p. 180]{Tohoku}
\end{reference}
In the situation of Lemma
\ref{lemma-special-coverings-compute-acyclicity}, suppose there is a
$d\geq0$ such that every $\mathcal R\in\mathfrak R(A)$ has nerve of
dimension at most $d$. If the equivalent conditions of that lemma hold for
$1\leq i\leq d$, then
$$
H^i(A,\mathcal F|_A)=0
\quad\text{for every }A\in\mathcal S\text{ and every }i>0.
$$
\end{lemma}

\begin{proof}
Choose a total ordering on the index set of $\mathcal R$. The usual
projection to alternating cochains and its simplicial homotopy inverse, as
in Lemma \ref{lemma-alternating-usual}, use only restriction maps and hence
apply to the subset covering considered here. The alternating complex is
zero in degrees greater than $d$. Consequently
$H^i(\mathcal R,\mathcal F)=0$ for $i>d$. By hypothesis it also vanishes
for $1\leq i\leq d$, so $A(m)$ holds for every $m$. The proof of Lemma
\ref{lemma-special-coverings-compute-acyclicity} shows that $B(m)$ holds for
every $m$. Applying Lemma \ref{lemma-covering-acyclicity-conditions} for
arbitrary $m$ proves the assertion.
\end{proof}
